% ============================================================================
\chapter{$\SL_2(\Z)$ 的模形式}
\label{ch:modular-forms-sl2z}
% ============================================================================

% ----------------------------------------------------------------------------
\section{模形式的定义}
% ----------------------------------------------------------------------------

\begin{definition}[模形式]
设 $k$ 为非负偶整数。$\SL_2(\Z)$ 上权为 $k$ 的\textbf{模形式}是一个全纯函数
$f\colon \HH \to \C$，满足：
\begin{enumerate}
  \item \textbf{变换律}：对任意
    $\gamma = \begin{pmatrix} a & b \\ c & d \end{pmatrix} \in \SL_2(\Z)$，
    \[
      f\!\left(\frac{a\tau + b}{c\tau + d}\right) = (c\tau + d)^k f(\tau);
    \]
  \item \textbf{在尖点全纯}：$f$ 在 $\infty$ 处全纯，即 Fourier 展开
    \[
      f(\tau) = \sum_{n=0}^{\infty} a_n q^n, \quad q = e^{2\pi i \tau},
    \]
    的系数 $a_n = 0$ 对所有 $n < 0$ 成立。
\end{enumerate}
权 $k$ 的模形式空间记作 $\Mk_k(\SL_2(\Z))$。
\end{definition}

\begin{definition}[尖点形式]
若进一步 $a_0 = 0$（即 $f$ 在尖点处消失），则称 $f$ 为\textbf{尖点形式}。
尖点形式空间记作 $\Sk_k(\SL_2(\Z))$。
\end{definition}

% ----------------------------------------------------------------------------
\section{$q$-展开原理}
% ----------------------------------------------------------------------------

% TODO: q-展开原理的陈述与证明

% ----------------------------------------------------------------------------
\section{模形式的例子}
% ----------------------------------------------------------------------------

% TODO: Eisenstein 级数 G_k，判别形式 Δ，j-不变量

% ----------------------------------------------------------------------------
\section{赋值公式（Valence Formula）}
% ----------------------------------------------------------------------------

% TODO: 赋值公式的陈述与详细证明

% ----------------------------------------------------------------------------
\section{习题}
% ----------------------------------------------------------------------------
